\section{Security Analysis and Type Checking Examples}

\subsection{Type Checking Examples}

\textbf{Example 1: Safe Program (Image Processing)}

\begin{verbatim}
Input:  medical_image : High
Output: result : ?

Program:  result := enhance(medical_image)

Type Check:
  enhance is type-preserving: High → High
  Assignment: High ⊑ result_type? 
  Yes, if result is declared High or inferred as High
  Verdict: ✓ SAFE (output labeled private/High)
\end{verbatim}

The output is private and only the owner can access it.

\textbf{Example 2: Explicit Leak (Malicious Program)}

\begin{verbatim}
Input:  secret : High, public_out : Low
Output: result : ?

Program:
  result := secret + 1
  if (result > 1000) then
    public_out := 1
  else
    public_out := 0

Type Check:
  Line 1: High + Low = High ⊔ Low = High
          result : High (inferred)
  Line 2: condition is High, so pc = High
  Line 4: Assignment public_out : Low
          Check: pc ⊔ 1 ⊑ Low?
          High ⊔ Low ⊑ Low? NO!
  Verdict: ✗ TYPE ERROR - Implicit leak detected
\end{verbatim}

The program is rejected. The malicious attempt to leak via control flow is blocked.

\textbf{Example 3: Implicit Leak via Timing}

\begin{verbatim}
Input:  secret : High, counter : Low
Output: counter

Program:
  while (secret > 0):
    counter := counter + 1
    secret := secret - 1
  return counter

Type Check:
  Condition: secret is High, so pc = High
  Assignment in loop: counter : Low
  Check: pc ⊔ (counter+1) ⊑ Low?
  High ⊔ Low ⊑ Low? NO!
  Verdict: ✗ TYPE ERROR - Timing leak detected
\end{verbatim}

The loop count would reveal the value of the secret. Type checking blocks this.

\subsection{Soundness Property}

\textbf{Theorem}: If a program type-checks with security types $\tau_1, \ldots, \tau_n$ for inputs and output type $\tau_{\text{out}}$, then:

\begin{equation}
\forall i: \tau_i \sqsubseteq \tau_{\text{out}} \text{ or input } i \text{ is unused}
\end{equation}

\textbf{Corollary}: High input cannot produce Low output. Malicious programs that attempt explicit or implicit leaks are rejected at type-check time.

\subsection{Hospital Application}

\textbf{Scenario}: Process medical images for enhancement and anonymize for research.

\begin{itemize}
    \item \textbf{Image Enhancement}: Input (private/High) → Output (private/High). Only owner accesses result.
    \item \textbf{Anonymization}: Input (private/High) → Output (public/Low) via approved declassifier. Researchers can access anonymized output.
    \item \textbf{Statistical Analysis}: Aggregate patient measurements → Produce population statistics with differential privacy noise. Output labeled Low.
\end{itemize}

In all cases, the Volpano type system ensures that no individual patient data leaks beyond what was explicitly approved via declassification.
